\scanpage{276}
\mspage{17}

Néanmoins on se gardera de \deleted{\illegible{1}} confondre un chemin \inserted{de type $I$} avec le chemin opposé de type $I^{\mathrm{opp}}$, la source de l'un \unsure{étant} le but de l'autre et inversement. C'est dire qu'un chemin doit être vu comme \uline{plus} qu'un \inserted{simple} foncteur
\diagram{\begin{tikzcd}[ampersand replacement=\&] I \arrow[r] \& X \end{tikzcd}},
il faut de plus se donner l'\uline{orientation de} $I$ comme faisant partie de la structure de chemin.

Donc la définition formelle raisonnable, qui colle à l'intuition \unsure{du chemin|des chemins} comme allant d'une \uline{source} bien définie vers un \uline{but} également bien défini, c'est qu'un chemin (défini en termes d'une cat. d'intervalle $I$) est la donnée d'un \uline{couple} $(\omega,c)$, où $\omega$ est une orientation de $I$, et \deleted{\unsure{$c$|,}}$c$ le diagramme $c\colon I\mathrel{\unsure{\text{---}|\longrightarrow}}X$.

Quand on regarde la catégorie des chemins d'un ``type'' déterminé, ce \unsure{type} est $(I,\omega)$, et non $I$. C'est quand $\omega$ est fixé \inserted{une bonne fois} \unsure{que} \deleted{la \unsure{ch}} cette catégorie s'identifie \unsure{simplement|impunément} à \deleted{\unsure{$\mathrm{Ch}$}} $\underline{\mathrm{Hom}}(I,X)$ sans plus.

Il faut voir cette catégorie de chemins comme munie de la structure essentielle
\diagram{\begin{tikzcd}[ampersand replacement=\&]
\underline{\mathrm{Ch}}_{(I,\omega)}\deleted{\unsure{(}}(X) \arrow[d,"{(\unsure{\sigma|\alpha},\beta)}"]\\
X\times X
\end{tikzcd}}
\[ \begin{aligned} c&\longmapsto(\sigma(c),\beta(c))\\[-2pt] &\text{où }\sigma(c)\text{ est la \uline{source},}\\ &\phantom{\text{où }}\beta(c)\text{ le but de }c. \end{aligned} \]
% archive sheet number: 275
